% =====================================================================
%  学習指導案：特別の教科 道徳（中学校）
%  ---------------------------------------------------------------------
%  単元：画面越しに友達と分かり合うこと
%  本時：AIとSNSを題材に、他者を理解することを考える
%  形式：オンライン非同期型授業
% =====================================================================
\documentclass[uplatex,a4paper,11pt]{jsarticle}

\usepackage{geometry}
\geometry{top=20mm,bottom=20mm,left=20mm,right=18mm}

\usepackage{array}
\usepackage{longtable}
\usepackage{enumitem}

\setlength{\tabcolsep}{4pt}
\renewcommand{\arraystretch}{1.35}

\newcommand{\itemhead}[2]{\par\medskip\noindent\textbf{#1.\quad #2}\par\nopagebreak\smallskip}

\pagestyle{plain}

\begin{document}

\begin{center}
  {\LARGE \textbf{特別の教科 道徳 学習指導案}}
\end{center}

\begin{flushright}
  \begin{tabular}{ll}
    指導教諭   & 〇〇　〇〇　　印 \\
    教育実習生 & 吉岡　優　　印 \\
  \end{tabular}
\end{flushright}

\bigskip

\noindent
1.\quad 日\hspace{1zw}時：令和8年5月26日（火）　第〇校時相当（オンライン非同期型授業／想定学習時間50分） \par
\noindent
2.\quad 実施学級：東京都〇〇区立　〇〇中学校　第2学年〇組　在籍生徒数　36名 \par
\noindent
3.\quad 単\hspace{1zw}元：特別の教科 道徳『中学道徳〇〇』（〇〇出版）　内容項目B「主として人との関わりに関すること」-(9)相互理解、寛容 \par


\itemhead{4}{単元観}
\hspace{1zw}本単元は、中学校学習指導要領（平成29年告示）第3章「特別の教科 道徳」第2「内容」B「主として人との関わりに関すること」(9)「相互理解、寛容」に基づいて設定したものである。同項目は「自分の考えや意見を相手に伝えるとともに、それぞれの個性や立場を尊重し、いろいろなものの見方や考え方があることを理解して、寛容の心をもって謙虚に他に学び、自らを高めていくこと」を求めている。\par
\hspace{1zw}「特別の教科 道徳」は平成29年告示の学習指導要領改訂において、「考え、議論する道徳」への質的転換が示された。本単元はその趣旨に基づき、生徒一人一人が多面的・多角的に思考し、自己を見つめる学びの場とすることを目指す。\par
\hspace{1zw}現代の中学生にとって、SNSやメッセージアプリは日常的なコミュニケーション手段であり、加えて生成AIとの対話も身近なものとなりつつある。これらの画面越しのやり取りの中で、「相手を理解する」ことはどのような意味を持つのか。本単元では、公開SNS・DM・AIキャラ・対面という4つの異なるコミュニケーション形態を対比的に扱い、「人と人が分かり合う」ことの本質を考えさせる。AIを「自己内省の鏡」として位置付け、AIにできることと人だからこそできることの違いを生徒自身に発見させたい。\par

\itemhead{5}{学級・生徒観}
\hspace{1zw}第2学年〇組は36名（男子19名・女子17名）。GIGAスクール構想により1人1台の学習用端末が配付されており、全員が自宅から学校LMSにアクセス可能な環境にある。\par
\hspace{1zw}SNS利用率は事前調査で約8割を超え、多くの生徒がメッセージアプリでの日常的なやり取りを行っている。一方で、既読プレッシャー・グループからの孤立感・他者の投稿を見て揺れる感情など、SNS特有の悩みを抱える生徒も少なくない。生成AIに触れた経験は約6割の生徒にあるが、AIとの対話を「相談」として行った経験はまだ限定的である。\par
\hspace{1zw}思春期特有の繊細さから、クラス全体としては比較的素直で穏やかであるが、自分の内面や感情を文章で表現することに抵抗を持つ生徒が一定数いる。特に、他者の前で自己開示することへのハードルが高い。本時では完全匿名性と「考えるだけ」モードの提供により、心理的安全性を確保する設計とした。\par

\itemhead{6}{指導観・教材観}
\hspace{1zw}本時の中心発問は「画面越しに、私たちは友達と分かり合えるのか」である。教材は、親友の公開SNS投稿が暗くなり、後日その親友から「最近きつい」とDMが届くという統合シナリオを採用する。これにより、公開SNSにおける傍観者問題と、DMにおける即時応答プレッシャーを同一文脈で扱う。\par
\hspace{1zw}教材内のAIキャラは、生徒の入力に応じて色が変化する「感情の鏡」として設計し、プルチックの感情の輪（1980）に基づく8色の色彩設計を採用する。AIキャラを擬人化しすぎず、あくまで「対話を通じて自己を見つめる装置」として位置付けることで、AIへの過度な依存を回避する。\par
\hspace{1zw}非同期型授業における対話的学びは、教師が事前に厳選した3つのサンプル振り返り（DM返答派／公開時の気づき派／メディア超越派）を読む活動、およびAIによる多面的視点の提示によって担保する。これは現代的なオンライン環境における「考え、議論する道徳」の実現を目指す試みである。\par
\hspace{1zw}本時はSNSとAIという題材の繊細さから、安全性設計を最重要事項として扱う。危険発話の検知・相談窓口の常時提示・教師アラートの3層構造により、生徒の心理的安全と緊急時対応を担保する。\par
\hspace{1zw}本単元は同テーマで設計された高等学校情報科の教材と「他者への想像力」という共通テーマで並列構造を取り、校種を超えた一貫した学びを目指す試みの一環である。\par

\itemhead{7}{単元の目標}
\begin{enumerate}[label=(\arabic*),leftmargin=2.8em,itemsep=3pt,topsep=2pt]
  \item AIとの対話・SNS・DM・対面など複数のコミュニケーション形態の特性と限界を理解する。（道徳的諸価値の理解）
  \item 親友との関係について多面的・多角的に考え、画面越しの言葉と直接の関わりの違いを自分なりに捉え直す。（物事を多面的・多角的に考える）
  \item 他者を理解しようとする態度を持ち、自分のコミュニケーションのあり方を見つめ直そうとする。（自己の生き方）
\end{enumerate}

\itemhead{8}{単元の指導計画（4時間扱い）}
\nopagebreak
{\renewcommand{\arraystretch}{1.4}
\begin{longtable}{|>{\centering\arraybackslash}m{16mm}|p{48mm}|p{96mm}|}
  \hline
  \multicolumn{1}{|c|}{時} & \multicolumn{1}{c|}{学習活動} &
  \multicolumn{1}{c|}{評価規準【観点】（評価方法）} \\ \hline
  \endfirsthead
  \hline
  \multicolumn{1}{|c|}{時} & \multicolumn{1}{c|}{学習活動} &
  \multicolumn{1}{c|}{評価規準【観点】（評価方法）} \\ \hline
  \endhead
  1 &
    人とのコミュニケーション（教科書教材）\newline pp.〇〇--〇〇 &
    ・対面・電話・文字などコミュニケーション形態の違いに気づいている。\newline 【道徳的諸価値の理解】（発言・ワークシート） \\ \hline
  2 &
    思いやりとは何か（教科書教材）\newline pp.〇〇--〇〇 &
    ・他者を思いやることの難しさと意味を多面的に考えている。\newline 【物事を多面的・多角的に考える】（記述） \\ \hline
  3\newline（本時） &
    画面越しに友達と分かり合えるのか（AIとSNSを題材として）\newline pp.〇〇--〇〇 &
    ・AIと人、テキストと対面の違いを多面的に考え、自己のコミュニケーションを見つめ直している。\newline 【多面的・多角的／自己の生き方】（振り返り記述） \\ \hline
  4 &
    他者の心を想像する（教科書教材）\newline pp.〇〇--〇〇 &
    ・他者理解の意義と方法を、本時の学びと結び付けて深めている。\newline 【自己の生き方】（記述） \\ \hline
\end{longtable}}

\itemhead{9}{本時の目標}
\begin{itemize}[label={○},leftmargin=2.2em,itemsep=3pt,topsep=2pt]
  \item AIとの対話・SNS・DM・対面という4つのコミュニケーション形態の特性と限界を理解する。【道徳的諸価値の理解】
  \item 親友との関わり方について多面的・多角的に考え、画面越しの言葉と人だからこそできる関わりの違いを自分なりに捉え直す。【物事を多面的・多角的に考える／自己の生き方】
\end{itemize}

\hspace{1zw}\textbf{※評価について：}道徳科の評価は数値による評定ではなく、個人内評価による文章評価を原則とする。本時においても、振り返り記述を中心に、生徒の多面的・多角的な思考の深まりと自己の生き方への接続を、個別に文章で評価する。\par

\itemhead{10}{本時の授業展開}
\nopagebreak
{\renewcommand{\arraystretch}{1.4}
\begin{longtable}{|>{\centering\arraybackslash}m{14mm}|p{72mm}|p{68mm}|}
  \hline
  \multicolumn{1}{|c|}{\shortstack{段階\\（時間）}} &
  \multicolumn{1}{c|}{学習活動（○指示・発問）} &
  \multicolumn{1}{c|}{指導上の留意点／☆評価の観点} \\ \hline
  \endfirsthead
  \hline
  \multicolumn{1}{|c|}{\shortstack{段階\\（時間）}} &
  \multicolumn{1}{c|}{学習活動（○指示・発問）} &
  \multicolumn{1}{c|}{指導上の留意点／☆評価の観点} \\ \hline
  \endhead
  導入\newline10分 &
    ・LMSにログインし、本時の教材ページにアクセスする。\newline
    ・イントロ動画（3分）を視聴する。\newline
    ○（基発問）「みなさんは普段、友達と何で連絡を取っていますか。」\newline
    ・普段使っているコミュニケーション手段を1〜2個答える。\newline
    ・教材ページに表示される橋渡し文（「SNS、メッセージアプリ、直接会って話す、いろいろあると思います。今日はその中でSNSとメッセージアプリ（DM）を題材に、友達と分かり合うことを考えていきます。」）を読む。\newline
    ・教材ページ上で、シチュエーション（公開SNSの暗い投稿、後日のDM）を2枚のスマホ画面モックで確認する（3分）。\newline
    ○（指示）「今のあなたの気持ちに一番近い色を、画面上の8色から1〜2色選んでタップしてください。正解はありません。今の自分の感じ方を大切にしてください。」 &
    ・教師は配信開始時刻に教師管理画面でログイン状況を確認し、未着手の生徒には10分後にリマインドを送信する。\newline
    ・基発問は本時の思考の外枠を決めるための導入問いとして、生徒の既有経験を引き出すことに集中させる。本時の核心となる問いは主発問（展開段階）で扱う。\newline
    ・基発問・主発問は、動画内で口頭で問いかけると同時に、教材ページにテキストで再掲して記入させる二重構造とする。\newline
    ・基発問への応答後、教材ページに橋渡し文を自動表示する設計とする。\newline
    ・教師制作のイントロ動画は3分以内とし、本時の流れを簡潔に示す。\newline
    ・シチュエーションはSNSの種類を特定せず「親友のSNS」と抽象化する。\newline
    ・色の選択は「正解」がないことを動画内で繰り返し明示し、自分の感じ方を大切にするよう促す。\newline
    ・自己開示への抵抗が予想される生徒に向け、「考えるだけ」モードが選択可能であることを最初の画面で説明する。 \\ \hline
  展開\newline30分 &
    ・教材内のAIキャラに状況を相談する（8分）。\newline
    ○（指示）「画面の中のAIキャラに、今の状況を話してみましょう。親友に何と返したいか、相談してみてください。書きたくない人は『考えるだけ』モードのままで構いません。」\newline
    ・AIキャラの色変化と応答を見て、自分の気持ちと比較する（2分）。\newline
    ・動画②「AIにできること、人だからこそできること」（6分）を視聴する。\newline
    ○（主発問）「AIキャラはあなたに、何らかの返事の案を提示してくれました。でも、それは本当にあなたが親友に送りたい言葉でしたか。AIに頼める言葉と、自分で書くしかない言葉。その違いは何だと思いますか。」\newline
    ・3つのサンプル振り返り（DM返答派／公開時の気づき派／メディア超越派）を読む（8分）。\newline
    ・AIから提示される多面的視点を確認する（3分）。\newline
    ○（中間問い）「3人の先輩たちは同じ状況で、それぞれ違う答えを書いています。あなたの考えに一番近いのはどれですか。また、自分にはなかった視点はどれですか。」&
    ・3つのサンプル振り返りは教師が事前に選定する。前年度生徒の記述から匿名化して採用するか、初回運用時は教師が複数の価値観を反映して作成する。\newline
    ・AIによる多面的視点は、生徒の応答内容に応じて自動生成する設計とする。\newline
    ・AIキャラの応答は断定を避け、複数の見方を示すようプロンプト設計してある。\newline
    ・AIキャラとの対話中に、危険発話（自殺念慮、自傷、攻撃性等）を検知した場合、AIは相談窓口（24時間子供SOSダイヤル0120-0-78310、よりそいホットライン0120-279-338、チャイルドライン0120-99-7777）を提示し、教師管理画面に即時アラートを送信する。\newline
    ・教師は管理画面で5分おきに各生徒の進捗を確認する。AIキャラとの対話が極端に短い／長い生徒には個別に状況を確認する。\newline
    ・サンプル振り返りは「正解」ではなく「多面的な考え方の例」として位置付け、「先輩たちはこう考えた。あなたはどう？」という構図で提示する。\newline
    ☆AIと人の違いに関する気づきを記述している。（AIキャラ対話内容・コメント記述） \\ \hline
  整理\newline10分 &
    ・多段階の振り返りを記入する（8分）。\newline
    ○（指示）「次の3つの問いに、自分の言葉で答えてください。長くなくて構いません。」\newline
    （1）公開SNSの投稿に気づいた時、あなたなら何をしましたか。\newline
    （2）DMが届いた今、あなたなら何を返したいですか。\newline
    （3）AIに頼める言葉と、自分で書くしかない言葉。その違いは何ですか。\newline
    ・本時のまとめ動画（2分）を視聴する。 &
    ・教師制作のまとめ動画では、「答えは出なくていい。これからも考え続けてほしい。」というメッセージで本時を締める設計とする。\newline
    ・振り返りは正解探しではなく、自分自身を見つめる時間として丁寧に促す。\newline
    ・記述内容は完全匿名で扱い、他生徒には個人の文章を公開しない設定とする。\newline
    ・教師は配信終了後、全生徒の振り返り記述を読み、個別文章評価コメントを48時間以内に返信する。\newline
    ・クラス全体の傾向（色彩分布・キーワード集計等）のみを匿名で次回授業で共有する。\newline
    ・参考文献（人間中心のAI社会原則、子供向けAIリテラシー資料、相談窓口等）へのリンクを提示する。\newline
    ☆多面的・多角的な思考の深まりが見られる。（振り返り記述）\newline
    ☆自己の生き方への接続が見られる。（振り返り記述） \\ \hline
\end{longtable}}

\bigskip

\noindent\textbf{【期待される振り返り記述の例】}

\hspace{1zw}\textbf{例1：}「私はAIキャラに『最近きついって言われた、どう返したらいい？』と聞いた。AIは『心配だね、そばにいるよと伝えてみたら』と優しい言葉を返してくれた。でも、それは『誰にでも当てはまる』言葉のような気がした。親友は、AIではなく『私』に向けてDMを送ってくれた。だから私が返す言葉も、『私だから書ける』ものでなければならないと思った。AIにできるのは『考えるきっかけ』を与えてくれることで、最終的に何を言うかは自分で決めなければならない。画面越しに分かり合うことは難しいけれど、文字の向こうに親友がいることを忘れずに、自分の言葉で向き合いたい。」

\hspace{1zw}\textbf{例2：}「正直に言うと、私は公開のSNS投稿が暗かった時、何もしなかった。誰か他の人がリアクションするだろうと思ったから。でも今振り返ると、私が気づいていたなら、私が動くべきだったのかもしれない。AIは『大丈夫？』と返すのを勧めたけど、私はDMよりも電話したい。文字では伝わらないものがあるから。AIは便利だけど、本当に大切な瞬間には、AIでは足りないと感じた。」

\bigskip

\noindent\textbf{【補足】指導上の特記事項}

\begin{itemize}[leftmargin=2em,itemsep=2pt,topsep=2pt]
  \item \textbf{非同期型授業特有の対応：}対面授業における机間巡視・板書・一斉発問は、それぞれ教師管理画面での進捗監視・教材ページのデザイン・動画内ナレーション及び教材内テキストに置き換えて設計している。教師は配信時間中、管理画面で全生徒の進捗を継続的に監視し、個別対応を行う。
  \item \textbf{心理的安全性の確保：}全演習に「考えるだけ」モードを提供し、自己開示を強制しない。記述内容は完全匿名で扱い、他生徒に個人の文章は公開しない。自分の入力はいつでも削除できる。
  \item \textbf{危機介入プロトコル：}AIキャラとの対話中に、深刻な悩み・自殺念慮・自傷・他者への攻撃などの危険発話を検知した場合、AIは相談窓口を提示し、教師管理画面に即時アラートを送信する。教師は24時間以内に生徒・保護者に連絡し、必要に応じてスクールカウンセラー・専門機関と連携する。
  \item \textbf{相談窓口（生徒に常時提示）：}24時間子供SOSダイヤル0120-0-78310／よりそいホットライン0120-279-338／チャイルドライン0120-99-7777
  \item \textbf{使用するリソース：}
    \begin{itemize}[leftmargin=1em,topsep=0pt,itemsep=0pt]
      \item 動画①「画面越しのコミュニケーション」（3分／教師制作）
      \item 動画②「AIにできること、人だからこそできること」（6分／教師制作）
      \item Webアプリ：AI感情対話アプリ（教師制作、Claude Haiku 4.5を内部API利用）
      \item AIキャラの色彩設計：プルチックの感情の輪（1980）に基づく8色
      \item 3つのサンプル振り返り：教師制作（前年度生徒振り返りより匿名化して選定）
    \end{itemize}
  \item \textbf{参考文献（生徒向け）：}
    \begin{itemize}[leftmargin=1em,topsep=0pt,itemsep=0pt]
      \item 内閣府「人間中心のAI社会原則」（\texttt{https://www8.cao.go.jp/cstp/aigensoku.pdf}）
      \item 文部科学省「初等中等教育段階における生成AIの利用に関するガイドライン」
      \item 法務省「子どもの人権110番」（\texttt{https://www.moj.go.jp/JINKEN/jinken112.html}）
    \end{itemize}
  \item \textbf{AIキャラの応答品質：}プロンプト設計において「中学生が相手であること」「断定しないこと」「専門的助言を避け信頼できる大人への相談を促すこと」「個人特定情報には言及しないこと」を明示し、定期的に応答品質をレビューする。
  \item \textbf{個人情報配慮：}入力された文章は暗号化保存し、保管期間（年度末まで）を明示する。実在の友人・家族の氏名等が含まれていた場合、AI側で自動マスキングする。
  \item \textbf{著作権配慮：}教材内の全素材は自作またはライセンスフリーであり、参考文献は公的機関の一次ソースに直リンクする。AI生成物は明示する。
  \item \textbf{教科横断的接続：}本教材は高校情報科「コミュニケーションと情報デザイン」の教材と共通テーマ「他者への想像力」で設計されており、校種を超えた一貫した学びを目指す試みの一環である。
\end{itemize}

\end{document}
